\documentclass[11pt,a4paper]{article}
\usepackage[T1]{fontenc}
\usepackage[utf8]{inputenc}
\usepackage{mathptmx}
\usepackage{amsmath,amssymb,amsthm}
\usepackage{amsfonts}
\usepackage{microtype}
\usepackage{geometry}
\geometry{margin=2.5cm}
\usepackage{natbib}
\usepackage{hyperref}
\usepackage{enumitem}
\usepackage{booktabs}
\usepackage{xcolor}

% ── Theorem environments ────────────────────────────────────────────────────
\newtheorem{theorem}{Theorem}[section]
\newtheorem{definition}[theorem]{Definition}
\newtheorem{proposition}[theorem]{Proposition}
\newtheorem{corollary}[theorem]{Corollary}
\newtheorem{claim}[theorem]{Claim}
\theoremstyle{remark}
\newtheorem{remark}[theorem]{Remark}

% ── Macros ──────────────────────────────────────────────────────────────────
\newcommand{\Pact}{\mathcal{P}_{\mathrm{act}}}
\newcommand{\Hon}{\mathcal{H}_{\mathrm{on}}}
\newcommand{\AP}{\mathcal{AP}}
\newcommand{\GM}{G_{M}}
\newcommand{\ARA}{A_{\mathrm{RA}}}
\newcommand{\MB}{\mathrm{MB}}
\newcommand{\GG}{\mathcal{G}}
\newcommand{\VV}{\mathcal{V}}
\newcommand{\EE}{\mathcal{E}}
\newcommand{\sigmaedge}{\sigma}

\begin{document}

% ── Title ────────────────────────────────────────────────────────────────────
\title{Algorithmic Causal Dynamics:\\
A Unified Framework for Emergence and Complexity\\
in Relational Actualism}

\author{Joshua F.\ Sandeman\\
\small Independent Researcher, Salem, Oregon\\
\small \texttt{sansuikyo@gmail.com}}
\date{March 2026}
\maketitle
\begin{center}
  {\small\textit{Preprint:}~\href{https://doi.org/10.5281/zenodo.19198171}
    {doi.org/10.5281/zenodo.19198171}\par}
\end{center}

% ── Abstract ─────────────────────────────────────────────────────────────────
\begin{abstract}
What does it mean for something to be more complex than something else? Standard
dynamical systems theory has no physical answer: complexity is a description
imposed on trajectories through a pre-existing phase space, not a property
generated by the dynamics themselves. This paper provides a physical answer
grounded in Relational Actualism (RA). We propose that the universe is a
discrete, growing causal Directed Acyclic Graph (DAG), and that every tier of
physical complexity---from fundamental particle interactions to biological cells
to minds---is a structural feature of this graph, generated by the same
underlying mechanism: the recursive coarse-graining of on-shell actualization
events into stable causal closures.

We formalize this via a hierarchy of recursive coarse-graining operators
$(\Pi_n)$ mapping microscopic actualization events to macroscopic Structural
Causal Models. This yields a precise topological definition of irreducible
emergence: a subgraph emerges at scale $n$ when it achieves sufficient causal
closure to function as a single effective node at scale $n+1$. Applying this
framework generates a four-tier continuum of physical complexity: (1) fundamental
interactions, (2) dissipative structures, (3) autopoietic systems, and (4) agentic
intelligence---each tier distinguished not by kind but by the depth and stability
of its causal closure.

The result is a unified physical account of complexity in which chemistry,
biology, and cognition are not separate domains requiring separate explanations,
but successive levels of the same topological process. We derive the mathematics
of each tier from first principles, grounding the Action Value function for
agentic graph-steering in the Steering Cost Functional of RACI
Theorem~5~\citep{Sandeman2026c} and establishing the formal correspondence
between the RA thermodynamic framework and Friston's Free Energy Principle.
New result: the Tier~2/3 Causal Firewall threshold $\sigma = \lambda\tau_d\ell^3 = 1$
is now proved via the directed Erd\H{o}s-R\'{e}nyi theorem applied to the
Poisson-CSG causal graph, resolving Hard Wall~1.
\end{abstract}

\noindent\textbf{Keywords:} Relational Actualism (RA) \textperiodcentered{} Causal DAG
\textperiodcentered{} Irreducible Emergence \textperiodcentered{} Coarse-Graining
\textperiodcentered{} Complexity \textperiodcentered{} Autopoiesis
\textperiodcentered{} Active Inference \textperiodcentered{} Graph Steering
\textperiodcentered{} Free Energy Principle

% ─────────────────────────────────────────────────────────────────────────────
\section{Introduction: Replacing the Phase Space with the Growing Graph}
\label{sec:intro}
% ─────────────────────────────────────────────────────────────────────────────

The question of what distinguishes a living cell from a crystal, or a mind from
a thermostat, is one of the deepest in science. Standard accounts appeal to
information processing, self-organization, or emergence---but these are
descriptions of complexity, not physical explanations of it. They characterize
what complex systems do; they do not identify the physical process that generates
the hierarchy of complexity in the first place, or explain why that hierarchy has
the specific tiers it does.

This paper offers a physical answer. Grounded in Relational Actualism
(RA)~\citep{Sandeman2026a}, we propose that every tier of physical complexity---
from fundamental particle interactions to biological cells to minds---is a
structural feature of the universal causal Directed Acyclic Graph (DAG)
$\GG = (\VV, \EE)$, generated by a single underlying mechanism: the recursive
coarse-graining of on-shell actualization events into stable causal closures.
Each vertex $v \in \VV$ is an irreversible on-shell actualization event---a
quantum transition in which a mediating boson crosses the kinematic threshold
$p^{\mu}p_{\mu} = m^2$ (natural units throughout) and an irreversible causal
fact is written into the structure of the universe. The macroscopic state space
is not a pre-existing manifold; it is the accumulated causal history of these
discrete events. The arena is continuously manufactured by the dynamics
occurring within it.

This shift has a precise consequence for the theory of complexity. A physical
system is not a point tracing a trajectory through an abstract phase space; it
\emph{is} a localized subgraph within the growing universal DAG. Its complexity
is determined by the depth and stability of its causal closure: how thoroughly
its internal dynamics are screened from its environment, and how many successive
levels of coarse-graining are required to represent it as a single effective node
at the next scale. Chemistry, biology, and cognition are not separate domains
requiring separate explanations. They are successive tiers of the same topological
process, each generated by the same formal mechanism applied at a higher level of
causal closure.

The present paper develops this picture into a formal framework we call
\emph{Algorithmic Causal Dynamics} (ACD). The central contribution is the
recursive coarse-graining hierarchy $(\Pi_n)$ and its consequences:

Section~\ref{sec:coarsegrain} formalizes the coarse-graining projection $\Pi_n$
mapping microscopic actualization events to macroscopic Structural Causal Models,
grounding it in the Markov blanket isomorphism of RACI~\citep{Sandeman2026c}.
Section~\ref{sec:hierarchy} establishes the recursive hierarchy and provides a
rigorous topological definition of irreducible emergence---the first such
definition grounded in causal graph physics rather than information-theoretic
postulates. Section~\ref{sec:continuum} derives the four-tier continuum of
physical complexity: fundamental interactions, dissipative structures, autopoietic
systems, and agentic intelligence, showing that each tier is distinguished by the
depth and stability of its causal closure rather than by any qualitative difference
in kind. Section~\ref{sec:steering} derives the mathematics of the highest tier
in full---the Action Value function for agentic graph-steering---directly from
the RA thermodynamic framework, unifying active inference with the Epistemic-Physical
Equivalence bound of RACI Theorem~5~\citep{Sandeman2026c} and establishing the
formal correspondence with Friston's Free Energy Principle. Sections~\ref{sec:predictions}
through~\ref{sec:hardwalls} state the empirical predictions, the relationship to
the RA suite, and the open problems.

% ─────────────────────────────────────────────────────────────────────────────
\section{Formalizing the Coarse-Grained Causal Graph}
\label{sec:coarsegrain}
% ─────────────────────────────────────────────────────────────────────────────

\subsection{The Microscopic Foundation}

Let the fundamental universe be described by the microscopic causal DAG
$\GG_0 = (\VV_0, \EE_0)$, where every vertex $v \in \VV_0$ is an irreversible
on-shell actualization event as defined in RAQM~\citep{Sandeman2026a}. The Local
Ledger Condition (RAGC, Definition~1~\citep{Sandeman2026b}) enforces
energy-momentum conservation at every vertex:
\begin{equation}
  \sum_{\substack{u \in \VV_0 \\ (u,v)\in\EE_0}} \mathbf{q}(u)
  \;=\;
  \sum_{\substack{w \in \VV_0 \\ (v,w)\in\EE_0}} \mathbf{q}(w),
\end{equation}
where $\mathbf{q}(v) \in \mathbb{R}^n$ is the charge vector (energy-momentum,
electric charge, baryon number, etc.) at $v$.

A macroscopic dynamical system does not compute every microscopic vertex. It
operates on a coarse-grained projection of the fundamental graph into a
macroscopic Structural Causal Model (SCM).

\begin{definition}[Macroscopic SCM]
A \emph{macroscopic Structural Causal Model} at scale $n$ is a DAG
$\GG_n = (\VV_n, \EE_n)$ together with:
\begin{itemize}[noitemsep]
  \item \emph{Exogenous variables} $U$: the underlying, untracked microscopic
    actualization events (thermal bath, environmental noise).
  \item \emph{Endogenous variables} $V$: the coarse-grained macroscopic states
    tracked at scale $n$.
  \item A structural equation $v_i = f_i(\mathrm{pa}(v_i), u_i)$ for each
    endogenous variable, where $\mathrm{pa}(v_i)$ are the causal parents of
    $v_i$ in $\GG_n$ and $u_i \in U$.
\end{itemize}
\end{definition}

\subsection{The Coarse-Graining Projection}

\begin{definition}[Coarse-Graining Projection]
The coarse-graining projection $\Pi_n : \GG_{n-1} \to \GG_n$ groups micro-histories
of $\GG_{n-1}$ into macroscopic equivalence classes. It is the repeated application,
at successively larger scales, of the Step~5 hydrodynamic limit of the Engine of
Becoming~\citep{Sandeman2026EB}: the same operation that produces the effective
spacetime metric from the actualized graph in RAGC produces the effective
biological and cognitive descriptions from the actualization graph in RAHC. The
four tiers of the complexity hierarchy are four scales of the same coarse-graining.

Two micro-histories $h_1, h_2 \in \GG_{n-1}$ are equivalent under $\Pi_n$ if and
only if they produce indistinguishable probability distributions over future
macroscopic observables at scale $n$:
\begin{equation}
  h_1 \sim_n h_2
  \;\iff\;
  P(\text{observables} \mid h_1) = P(\text{observables} \mid h_2).
  \label{eq:equiv}
\end{equation}
The vertices $\VV_n$ are the equivalence classes under $\sim_n$. The edges
$\EE_n$ represent the aggregate causal influence between these coarse-grained states,
induced by the causal structure of $\GG_{n-1}$.
\end{definition}

The equivalence relation~\eqref{eq:equiv} is the formal content of the intuition
that macroscopic descriptions discard microscopic detail that is irrelevant to
future macro-observables. It is the RA analogue of the coarse-graining implicit in
thermodynamics, but here applied to the causal graph directly rather than to a
phase space distribution.

\subsection{Connection to the Markov Blanket}

The coarse-graining projection $\Pi_n$ is not arbitrary: it is constrained by the
causal structure of $\GG_{n-1}$. Specifically, the boundary of the macroscopic
equivalence class $[h]_n$ is determined by the Markov blanket of the subgraph
being coarse-grained. By RACI Theorem~2 (Markov Blanket Isomorphism,
\citep{Sandeman2026c}), the Markov blanket $\MB(G_O)$ of an organism subgraph
$G_O \subseteq \GG$ is formally isomorphic to the boundary $\partial J^-(G_O)$
of its past light cone:
\begin{equation}
  X_{G_O} \;\perp\; X_{\VV \setminus (G_O \cup \partial J^-(G_O))}
  \;\Big|\; X_{\partial J^-(G_O)}.
\end{equation}
This means that the coarse-graining boundary is not a matter of convention or
epistemic choice; it is determined by the causal topology of the underlying graph.
The macro-node in $\GG_n$ that represents the coarse-grained system is precisely
the node whose incoming and outgoing edges correspond to the system's Markov
blanket.

% ─────────────────────────────────────────────────────────────────────────────
\section{The Recursive Coarse-Graining Hierarchy and the Topology of Emergence}
\label{sec:hierarchy}
% ─────────────────────────────────────────────────────────────────────────────

\subsection{Irreducible Emergence}

In standard physics, ``emergence'' is often treated as an epistemic phenomenon:
a property is emergent if it cannot be deduced from a tractable description of
the parts. This epistemic account makes emergence a matter of computational
convenience rather than physical reality.

Algorithmic Causal Dynamics rejects this. Because the universe is a growing graph
and because the coarse-graining projection $\Pi_n$ is constrained by the causal
topology of the underlying structure, emergence has a topological definition.

\begin{definition}[Irreducible Emergence]
A subgraph $S \subseteq \GG_{n-1}$ exhibits \emph{irreducible emergence} at
scale $n$ if:
\begin{enumerate}[noitemsep,label=(\roman*)]
  \item $S$ achieves sufficient causal closure that its Markov blanket
    $\MB(S)$ is stable under the growth dynamics of $\GG_{n-1}$; and
  \item The macroscopic edges $\EE_n$ incident on the coarse-grained node
    $\Pi_n(S) \in \GG_n$ cannot be recovered from any strict sub-decomposition
    of $S$ at scale $n-1$.
\end{enumerate}
\end{definition}

Condition (i) is the stability condition: the Markov blanket persists under
continued actualization of the environment. Condition (ii) is the irreducibility
condition: the causal influences that the emergent system exerts and receives at
the macro scale cannot be recovered by examining its parts separately. This is
not epistemic limitation; it is structural. The macroscopic edges $\EE_n$ are
real causal relationships that exist \emph{only} at the coarse-grained level
because they depend on the collective organization of the subgraph, not on any
individual vertex.

\subsection{The Recursive Hierarchy $(\Pi_1, \Pi_2, \ldots, \Pi_n)$}

Applying the coarse-graining projection recursively generates a nested hierarchy
of causal descriptions:

\begin{description}[leftmargin=1.5em]
  \item[\textbf{Level 0} ($\GG_0$)] The fundamental actualization events: on-shell
    boson exchanges satisfying $p^{\mu}p_{\mu} = m^2$.

  \item[\textbf{Level 1} ($\Pi_1 \to \GG_1$, \textit{Chemistry})] A localized
    sub-network of Level 0 events forms a stable covalent structure. To the
    surrounding environment, the internal actualization history is screened off
    by the molecule's Markov blanket. The molecule acts as a single causal
    node in $\GG_1$. This is the level at which the Assembly Index
    $A(M)$~\citep{Murray2023} becomes the relevant complexity measure, bounded
    below by $\ARA(M)$ per RACI Theorem~1~\citep{Sandeman2026c}.

  \item[\textbf{Level 2} ($\Pi_2 \to \GG_2$, \textit{Cells})] A massive network
    of Level 1 molecular nodes achieves autopoietic closure: a Markov blanket
    that actively maintains itself against environmental perturbation. To the
    surrounding tissue, the internal metabolic actualization history is
    irrelevant; the cell acts as a single effective node in $\GG_2$.

  \item[\textbf{Level 3} ($\Pi_3 \to \GG_3$, \textit{Organisms})] A network of
    cells achieves organismal-level Markov closure. The organism acts as a node
    in $\GG_3$ that interacts with other organisms and environmental structures.

  \item[\textbf{Level 4} ($\Pi_4 \to \GG_4$, \textit{Agentic Intelligence})]
    An organism develops a \emph{generative motif}---a persistent internal
    sub-DAG that is a homomorphic image of the environmental causal structure
    and runs predictions faster than the environment actualizes
    (Section~\ref{subsec:generative}). This is the level at which active
    inference~\citep{Friston2010} and graph steering (Section~\ref{sec:steering})
    operate.
\end{description}

\begin{remark}
The hierarchy is not a partition of natural kinds but a sequence of stable
causal closure regimes. The relevant level for describing a system is determined
by the scale at which its Markov blanket is stable. A system can participate in
multiple levels simultaneously: a neuron is a Level 2 node (cell) that is also
a component of a Level 3 node (organism) that steers at Level 4.
\end{remark}

\subsection{The End of Epistemic Reductionism}

The topological definition of irreducible emergence has a sharp consequence for
reductionism. Reducing a Level 4 agent to Level 0 atoms does not merely make the
description intractable; it destroys the macroscopic edges $\EE_4$ that govern
the system's actual causal behavior. The agentic steering described in
Section~\ref{sec:steering} operates exclusively over the macroscopic topology
$\GG_4$. No description at Level 0 or Level 1 recovers these edges, because they
depend on the collective organization of the full subgraph at Level 4.

This is not a claim about computational complexity. It is a claim about causal
structure: the edges $\EE_n$ for $n \geq 2$ are physically real causal
relationships that do not exist at lower scales. They are produced by the
coarse-graining process, which is itself constrained by the causal topology of
the underlying graph. Reductionism fails not because the details are too numerous
to track, but because the relevant causal structure is not present at the
microscopic level.

% ─────────────────────────────────────────────────────────────────────────────
\section{The Four-Tiered Continuum of Complexity}
\label{sec:continuum}
% ─────────────────────────────────────────────────────────────────────────────

The recursive coarse-graining hierarchy $(\Pi_n)$ is not a taxonomy imposed on
nature from the outside; it is a physical structure generated by the causal
topology of the universal DAG. Each tier in the following continuum corresponds
to a specific regime of causal closure, characterised by the stability and
depth of the Markov blanket achieved at that scale. The tiers are not
discontinuous kinds but thresholds along a single continuous process---the
progressive deepening of causal self-reference within the growing graph.

\subsection{Tier 1: Fundamental Interactions (The Microscopic Engine)}

At the base level $\GG_0$, the universe expands via local, irreversible
actualizations governed by the RA-CSG dynamics of RAGC~\citep{Sandeman2026b}.
When a mediating boson crosses the kinematic threshold $p^{\mu}p_{\mu} = m^2$,
an irreversible vertex is inscribed. There is no coarse-graining here; this is
the raw, uncompressed generation of causal history. The Local Ledger Condition
enforces conservation at every vertex; the LLC is the discrete substrate from
which all macroscopic conservation laws emerge in the continuum limit.

\subsection{Tier 2: Dissipative Structures (Reactive Macroscopics)}

Macroscopic dissipative structures---Rayleigh-B\'{e}nard convection cells,
hurricanes, chemical oscillators---emerge when a subgraph of $\GG_1$ provides an
efficient topological pathway for entropy production across a thermodynamic
gradient. They persist in time and possess a definable macroscopic SCM at Level 1.

The key distinction from autopoietic systems is the stability of the Markov
blanket. A dissipative structure's ``boundary'' is a statistical pattern in
$\GG_1$, but it is not maintained by internal steering: the subgraph has no
mechanism for detecting and correcting deviations from the boundary condition.
When the thermodynamic gradient vanishes, the subgraph dissolves seamlessly into
environmental noise. The criticality parameter $\sigmaedge$---the ratio of
internal actualization rate to environmental disruption rate~\citep{Sandeman2026c}---
is non-zero but unregulated: it is set entirely by external conditions rather than
by internal causal forcing.

\begin{remark}
Some dissipative structures (notably certain chemical oscillators) exhibit
rudimentary self-repair under perturbation. The Tier 2 / Tier 3 boundary is
therefore a continuum, not a sharp phase transition. The Causal Firewall conjecture
of RACI~\citep{Sandeman2026c} identifies the critical threshold
$\lambda \cdot \tau_d \cdot \ell^3 = 1$ as the candidate for a sharp boundary,
pending the limit theorem discussed in that paper's hard wall.
\end{remark}

\subsection{Tier 3: Autopoietic Systems (The Edge of Chaos)}

The emergence of biological life marks a phase transition in graph dynamics. An
autopoietic system~\citep{Maturana1980} achieves a \emph{stable Markov blanket}:
a causal boundary $\partial J^-(G_O)$ that is actively maintained against
environmental perturbation. By RACI Theorem~2~\citep{Sandeman2026c}, this
boundary is formally isomorphic to the Markov blanket of the organism in the
Bayesian network sense.

The defining property of Tier 3 is not merely the presence of a Markov blanket
but the active regulation of the criticality parameter $\sigmaedge$ toward the
edge of chaos ($\sigmaedge \approx 1$). An autopoietic system continuously executes
highly specific actualization sequences to maintain this regime:
\begin{itemize}[noitemsep]
  \item If $\sigmaedge < 1$ (subcritical): internal actualization rate is insufficient
    to maintain the graph structure against environmental noise. The Markov blanket
    dissolves (death).
  \item If $\sigmaedge > 1$ (supercritical): internal actualization rate exceeds
    the coherence threshold; the system loses the structured causal closure that
    distinguishes it from its environment (runaway growth, loss of functional
    identity).
  \item At $\sigmaedge \approx 1$ (critical): the system maintains maximum
    computational and causal complexity compatible with stable Markov closure.
\end{itemize}

This maintenance is thermodynamically costly. By RACI Theorem~4 (Landauer Cost
of Actualization~\citep{Sandeman2026c}), each actualization event that suppresses
$\Delta q$ bits of quantum uncertainty dissipates at minimum
$W_{\mathrm{act}} \geq \Delta q \cdot k_BT \ln 2$. Autopoietic persistence
requires a continuous metabolic flux to pay this cost.

\subsection{Tier 4: Agentic Systems (Active Inference and Graph Steering)}

Intelligence represents the highest topological survival strategy available to a
Tier 3 system. An agentic system extends its causal influence outward beyond its
Markov blanket by maintaining an internal generative model of the external
coarse-grained graph $\GG_{\mathrm{env}}$ and using that model to steer its own
future actualization sequence.

\subsubsection{The Topological Criterion for Agency: The Generative Motif}
\label{subsec:generative}

The formal distinction between Tier 3 and Tier 4 requires a precise topological
criterion. A Tier 3 autopoietic system maintains its Markov blanket reactively:
it detects perturbations and responds to them after they arrive. A Tier 4 agentic
system predicts perturbations before they arrive by running physical simulations
of the environment's future causal structure.

\begin{definition}[Generative Motif]
A subgraph $\GG_{\mathrm{model}} \subseteq G_O$ is a \emph{generative motif} if:
\begin{enumerate}[noitemsep, label=(\roman*)]
  \item $\GG_{\mathrm{model}}$ is a homomorphic image of $\GG_{\mathrm{env}}$ at
    the relevant coarse-grained scale: there exists a structure-preserving map
    $\phi: \GG_{\mathrm{env}} \to \GG_{\mathrm{model}}$ that preserves the
    causal partial order and the macroscopic edge structure $\EE_n$;
  \item The causal depth of $\GG_{\mathrm{model}}$ is sufficient to compute the
    probability distribution over future states of $\GG_{\mathrm{env}}$ before
    those states are actualized in the external graph; and
  \item The simulation is executed via real, on-shell actualization events within
    $G_O$---the generative motif is a physical sub-process, not a disembodied
    computational abstraction.
\end{enumerate}
\end{definition}

Condition (iii) is crucial and requires unpacking.

\subsubsection{The Thermodynamic Cost of Imagination}
\label{subsec:imagination_cost}

The generative motif is not floating in the off-shell potentia. Running a physical
model of the environment requires real, on-shell actualization events within the
agent's own graph --- synaptic firings, ATP hydrolysis, ion channel dynamics in
the biological case. These are irreversible causal events that write permanent
vertices to the agent's internal DAG. This has a direct thermodynamic consequence:

\begin{proposition}[Imagination Has a Landauer Floor]
\label{prop:imagination}
Each internal simulation step of the generative motif --- each counterfactual
evaluation of a future environmental state --- is a forced actualization sequence
$\pi_{\mathrm{sim}}$ that dissipates at minimum
\begin{equation}
  W_{\mathrm{steer}}(\pi_{\mathrm{sim}}) \;\geq\; k_BT \cdot \Delta F[q_{\mathrm{sim}}],
  \label{eq:imagination_cost}
\end{equation}
where $\Delta F[q_{\mathrm{sim}}]$ is the reduction in the agent's uncertainty
about the future environmental state achieved by the simulation step.
\end{proposition}

\begin{proof}
Direct application of the Epistemic-Physical Equivalence bound
(RAHC Theorem~5, Section~\ref{sec:steering_rahc}) to the internal simulation
sequence $\pi_{\mathrm{sim}}$.
\end{proof}

Equation~\eqref{eq:imagination_cost} has a direct empirical anchor: the human
brain's consumption of approximately 20\% of the body's resting metabolic
energy~\citep{Attwell2001} is a physical consequence of maintaining a
high-fidelity generative motif of the social and physical environment. This is
not an evolutionary inefficiency; it is the Landauer floor for running a
sufficiently detailed physical simulation of a complex world. The metabolic cost
of intelligence scales with the fidelity and breadth of the internal model, not
with the complexity of the actions it selects.

A further consequence: the agent cannot simulate arbitrarily many counterfactual
futures simultaneously. The Action Value function~\eqref{eq:actionvalue}
incorporates $W_{\mathrm{steer}}(\pi)$ as the cost of both external actions
\emph{and} internal simulations that evaluate them. The finite metabolic budget
of any physical agent therefore enforces a genuine trade-off between the depth
of predictive simulation and the breadth of policy search --- a thermodynamically
grounded account of the bounded rationality observed in all biological
intelligence~\citep{Kauffman2000}.

\begin{definition}[Tier 4 Transition Criterion]
A system crosses the $\Pi_4$ threshold from Tier 3 to Tier 4 if and only if
it contains a generative motif $\GG_{\mathrm{model}}$ whose internal recurrent
causal structure (macroscopic edges $\EE_3$ with sufficient loop density) can
compute predictions of $\GG_{\mathrm{env}}$ faster than $\GG_{\mathrm{env}}$
actualizes---i.e., the internal simulation rate exceeds the environmental
actualization rate at the relevant coarse-grained scale.
\end{definition}

The recurrence requirement---sufficient loop density in $\EE_3$---is the topological
signature of predictive capacity: recurrent causal loops enable the generative
motif to propagate predictions forward in simulated time. A purely feedforward
subgraph (no recurrent loops) can match the environment's current state but cannot
anticipate its future states.

\subsubsection{The Epistemic Nature of the Generative Motif's Content}
\label{subsec:epistemic}

A critical clarification is required to prevent a category error that theoretical
neuroscience reviewers will correctly identify. The Generative Motif definition
involves two distinct aspects that must be kept separate:

\textbf{The substrate is ontic.} The physical sub-DAG $\GG_{\mathrm{model}}$ is
a real causal structure within the agent's own graph. Its actualization events
--- synaptic firings, ATP hydrolysis, ion channel dynamics --- are irreversible
on-shell events that pay the Landauer thermodynamic cost. The generative motif
is not a mathematical abstraction; it is a physical sub-process. This is the
sense in which condition~(iii) of the Generative Motif definition is non-negotiable:
the substrate is as ontic as any other part of the causal DAG.

\textbf{The content is epistemic.} The \emph{content} of $\GG_{\mathrm{model}}$
--- the causal relationships it represents about the environment --- is an
epistemic approximation, not an ontological replica. The homomorphism
$\phi: \GG_{\mathrm{env}} \to \GG_{\mathrm{model}}$ in condition~(i) operates
at the coarse-grained macroscopic scale; it is a functional, lossy map, not a
micro-physical isomorphism. The model is not the territory.

\textbf{The ``good enough'' principle.} Evolution does not select for truth; it
selects for survival. The generative motif that persists is the one whose
coarse-grained predictions are adequate to prevent Markov blanket breach --- not
the one that most faithfully represents the microscopic physics of the environment.
A formally adequate generative motif need only satisfy:
\begin{enumerate}[noitemsep, label=(\roman*)]
  \item Causal-order preservation at the macroscopic scale: predictions respect
    the temporal structure of environmental actualization at the scale relevant
    to the agent's actions.
  \item Sufficient fidelity for action selection: the probability distribution
    over future environmental states computed by $\GG_{\mathrm{model}}$ must be
    accurate enough to distinguish between actions that preserve versus threaten
    Markov closure.
\end{enumerate}
A model that satisfies these two conditions is evolutionarily adequate regardless
of whether it faithfully represents microscopic causal structure. This is why
biological organisms maintain highly compressed, lossy world-models that are
systematically wrong about physics at the quantum level yet remain extraordinarily
effective at the ecological scale. The framework does not require biological
organisms to be physicists; it requires them to be adequate predictors of
ecologically relevant causal structure.

This distinction resolves a potential confusion about condition~(i) of the
Generative Motif definition. The homomorphism $\phi$ is not a claim that the
organism's internal model achieves micro-physical accuracy. It is a claim about
functional correspondence at the coarse-grained level: the macroscopic causal
edges that govern survival-relevant environmental events are structurally
represented in $\GG_{\mathrm{model}}$, even if the microscopic detail is massively
compressed or absent. The physical substrate paying the Landauer cost is ontic;
the representation it implements is epistemic.

This is the causal DAG formulation of active inference~\citep{Friston2010}:
the agent continuously minimizes the variational free energy between its internal
generative motif and its sensory evidence by (i) updating the motif (perception)
and (ii) selecting actions that steer the global DAG toward preferred states (action).

Section~\ref{sec:steering} derives the formal mathematics of graph steering from
the RA framework.



\subsubsection{Worked Example: The Three Levels in a Single Biological System}
\label{subsec:worked_example}

To prevent the category error identified above from arising in practice, we
trace a single biological system---\textit{E.\ coli} chemotaxis---through
all three levels of the RAHC architecture simultaneously. This illustrates
precisely how microscale on-shell events, coarse-grained causal nodes, and
epistemic model content are related without being conflated.

\paragraph{Level 0 / Tier 1: The on-shell substrate.}
At the microscale, the physical reality of the chemotaxis system consists
entirely of on-shell actualization events in the sense of
RAQM~\citep{Sandeman2026a}:
\begin{itemize}[noitemsep]
  \item \textbf{Ligand binding:} a maltose molecule undergoes an on-shell
  electromagnetic exchange with the MBP receptor binding pocket. This
  constitutes an irreversible actualization event that writes a vertex into
  the causal DAG and produces an environmental entanglement record.
  \item \textbf{Phosphorylation cascade:} CheA autophosphorylation and
  CheY-P production proceed via on-shell ATP hydrolysis events. Each
  hydrolysis event is a Kinematic Snap: a mediating boson (the
  $\gamma$-phosphate) crosses the mass shell, irreversibly transferring
  causal information down the signalling chain.
  \item \textbf{Flagellar motor switching:} CheY-P binding to the FliM
  switch complex triggers a conformational change via on-shell
  protein-protein interaction events, reversing flagellar rotation.
\end{itemize}
These are not metaphors. They are specific, catalogued on-shell events whose
molecular mechanisms are known to $\sim$\AA ngstr\"{o}m resolution. The
substrate is ontic.

\paragraph{Level 1 / Tier 2--3: The coarse-grained causal node.}
The coarse-graining projection $\Pi_2$ groups the $\sim 10^5$ individual
molecular actualization events per second into a single macro-node in
$\GG_2$: the \emph{CheY-P gradient}. This node is not a summary statistic
imposed by the observer; its boundary is fixed by the Markov blanket
constraint (RAHC Theorem 3): the CheY-P gradient is the minimal
coarse-grained state sufficient to predict the bacterium's future
motor-switching behaviour, given the current environmental sugar gradient.

At $\Pi_3$, the bacterium's motility pattern (run-and-tumble) is a
macro-node in $\GG_3$. The causal edge from the sugar gradient to the
motility pattern exists at the Tier 3 level and is \emph{not present}
in the Tier 0 description. It emerges from the collective organisation of
the signalling subgraph---this is the irreducible emergence of
Section~\ref{sec:hierarchy}.

\paragraph{Level 2 / Tier 3--4: The epistemic content.}
The bacterium does not have a Generative Motif in the full Tier 4 sense
(it lacks the recursive predictive model of an agentic system). But it has
a \emph{proto-Generative Motif}: the CheY-P gradient constitutes a
compressed, lossy causal representation of the external sugar gradient.
This representation is:
\begin{itemize}[noitemsep]
  \item \textbf{Ontic substrate:} real phosphorylation events in a real
  sub-DAG, paying a real Landauer cost.
  \item \textbf{Epistemic content:} a functional, lossy map of the external
  environment, adequate to drive up-gradient motility without faithfully
  representing the microscopic distribution of individual sugar molecules.
\end{itemize}

\paragraph{The architectural summary.}
The reviewer who asks ``how does the neural network implement graph-steering
in terms of on-shell events?'' is conflating Level 2 (epistemic content)
with Level 0 (on-shell substrate). The answer is:
\begin{enumerate}[noitemsep]
  \item \textbf{The substrate} of every biological computation consists of
  on-shell events (ATP hydrolysis, ion channel openings, photon absorption,
  synaptic vesicle release). These are Tier 1 actualisations. Named
  explicitly in Section~\ref{subsec:epistemic}.
  \item \textbf{The coarse-grained causal structure} at Tier 3--4 is
  real and causally efficacious but is not recoverable from the Tier 0
  description alone. It requires the $\Pi_n$ projection constrained by
  the Markov blanket.
  \item \textbf{The epistemic content} of the internal model is a lossy
  functional map at the macroscopic tier. Demanding that it be specified
  in terms of Tier 0 on-shell events is a category error equivalent to
  demanding that Newton's gravitational law be expressed in terms of
  individual quark-gluon interactions.
\end{enumerate}
The framework is not silent on the microscale implementation; it is
precise about which questions belong to which level of description.

% ─────────────────────────────────────────────────────────────────────────────
\section{Foundational Theorems of Causal Complexity}
\label{sec:theorems}
% ─────────────────────────────────────────────────────────────────────────────

The four-tier continuum of Section~\ref{sec:continuum} is not merely a
descriptive taxonomy. Each tier transition is governed by a precise mathematical
condition, and each condition yields a theorem. This section establishes the five
foundational theorems of causal complexity: the Assembly Index Correspondence
(Tier 1/2 boundary), the DAG Markov Property and Markov Blanket Isomorphism
(Tier 2/3 boundary), and the Landauer Cost and Epistemic-Physical Equivalence
(Tier 3/4 boundary). These results are proved here; their biological applications
are developed in the companion paper RACI~\citep{Sandeman2026c}.

% ── Tier 1/2: Assembly Index ─────────────────────────────────────────────────
\subsection{The Assembly Index Correspondence (Tier 1/2 Boundary)}
\label{sec:assembly_rahc}

The transition from Level 0 (fundamental interactions) to Level 1 (stable
molecular structures) is governed by the irreversibility of covalent bond
formation. The Assembly Theory framework of Murray et al.~\citep{Murray2023}
quantifies molecular complexity by the minimum number of JOIN operations in a
construction program. ACD provides the physical underpinning of this metric.

\begin{definition}[Actualism Assembly Index]
For a molecule $M$ assembled in a region of the causal DAG $\GG$, define the
\emph{molecular subgraph} $\GM$ as the induced subgraph on the set of
actualization events involved in the assembly of $M$. The Actualism Assembly
Index is
\begin{equation}
  \ARA(M) \;:=\; \mathrm{depth}(\GM)
  \;=\; \text{length of the longest directed path in } \GM.
  \label{eq:ara}
\end{equation}
\end{definition}

\begin{theorem}[Assembly Index Correspondence]
\label{thm:assembly}
Under Condition~C1$'$ (each JOIN operation corresponds to at least one irreversible
actualization event, operationally defined as an entropy-producing on-shell energy
transfer not reversed by subsequent unitary evolution) and Condition~C2 (the
ordering of operations is consistent with the causal partial order of $\GM$):
\begin{equation}
  \ARA(M) \;\geq\; A(M),
\end{equation}
with equality when $\GM$ is a minimal assembly DAG.
\end{theorem}

\begin{proof}
The Assembly Index $A(M)$ counts the minimum number of JOIN operations. Each
JOIN corresponds (by C1$'$) to at least one vertex in $\GM$. The minimum number of
sequential such vertices, respecting the causal partial order (C2), is the depth
of $\GM$. In a minimal DAG, depth equals the number of edges on the longest path
$= A(M)$. In a non-minimal DAG (with redundant parallel paths), depth $\geq A(M)$
because parallel paths do not reduce the sequential minimum.
\end{proof}

\begin{remark}
Condition C1$'$ is established by a variational argument (Born-Oppenheimer
approximation guarantees all stable covalent bonds are exothermic, mandating
irreversible energy dissipation) and confirmed computationally by a B3LYP/6-311+G*
density functional survey across the five primary bond types of biological assembly
(C-C, C-N, C-O, S-S, P-O). The full treatment, including the polymer-regime
coarse-graining rule and the linear scaling $\ARA(M) \approx \kappa \cdot A(M)$
in condensed phases, is given in RACI~\citep{Sandeman2026c}.
\end{remark}

% ── Tier 2/3: Markov Blanket ─────────────────────────────────────────────────
\subsection{The Markov Blanket Isomorphism (Tier 2/3 Boundary)}
\label{sec:markov_rahc}

The transition from Level 1 (chemistry) to Level 2 (autopoietic cells) requires
a precise topological criterion for causal closure. The Markov blanket isomorphism
establishes that this criterion is the equivalence between the organism's
statistical boundary (the Markov blanket) and its physical boundary (the past
light cone boundary in the causal DAG).

In a Bayesian network $N = (X, E_N)$, the Markov blanket $\MB(S)$ of a node set
$S$ is the minimal set rendering $S$ conditionally independent of all other nodes:
\begin{equation}
  X_S \perp X_{V \setminus (S \cup \MB(S))} \mid X_{\MB(S)}.
\end{equation}
In the RA causal DAG, define the past light cone
$J^-(v) := \{u \in V : u \text{ causally precedes } v\} \cup \{v\}$ and its
boundary $\partial J^-(v) := \{u \in J^-(v) : \exists\, w \notin J^-(v) \text{
with } (u,w) \in E\}$.

\begin{theorem}[DAG Markov Property]
\label{thm:dag-markov}
In the causal DAG $\GG$ with the RA-CSG probability measure, for any vertex $v$:
\begin{equation}
  \rho_v \;\perp\; \rho_{V \setminus J^-(v)} \;\Big|\; \rho_{\partial J^-(v)}.
\end{equation}
\end{theorem}

\begin{proof}
The state at $v$ is determined entirely by its causal predecessors $J^-(v)$.
The boundary $\partial J^-(v)$ screens off all further causal information: any
state $u \notin J^-(v)$ can influence $v$ only through paths passing through
$\partial J^-(v)$. The conditional independence follows from the d-separation
criterion on DAGs~\citep{Pearl1988}: $v$ is d-separated from $V \setminus J^-(v)$
by $\partial J^-(v)$.
\end{proof}

\begin{theorem}[Markov Blanket $\leftrightarrow$ Light Cone Boundary Isomorphism]
\label{thm:iso}
\textup{[Lean-verified core.]}
Let $O$ be an organism modelled as a subgraph $G_O \subseteq \GG$. There exists a
coarse-graining map $\varphi: \partial J^-(G_O) \to \MB(G_O)$ such that:
\begin{enumerate}[label=\emph{(\alph*)},noitemsep]
  \item Sensory states correspond to incoming boundary vertices of $\partial J^-(G_O)$.
  \item Active states correspond to outgoing boundary vertices of $G_O$.
  \item The conditional independence structure is preserved:
  \begin{equation}
    X_{G_O} \;\perp\; X_{V \setminus (G_O \cup \partial J^-(G_O))}
    \;\Big|\; X_{{\partial J^-(G_O)}}.
  \end{equation}
\end{enumerate}
\end{theorem}

\begin{proof}
Parts (a) and (b) are definitional correspondences. For (c): by
Theorem~\ref{thm:dag-markov}, $\rho_v$ is conditionally independent of
$V \setminus J^-(v)$ given $\partial J^-(v)$ for each $v \in G_O$. Since
$\partial J^-(G_O) \subseteq \partial J^-(v)$ for all $v \in G_O$, joint
conditional independence follows by the intersection property of conditional
independence on DAGs.
\end{proof}

% ── Tier 3/4: Landauer and Steering ─────────────────────────────────────────
\subsection{The Thermodynamics of Causal Forcing (Tier 3/4 Boundary)}
\label{sec:steering_rahc}

The transition from Level 3 (autopoietic organisms) to Level 4 (agentic
intelligence) requires that the system not merely maintain causal closure but
actively force the growth of the DAG toward preferred states. This forcing carries
a thermodynamic cost that bounds all intelligent behaviour.

\begin{theorem}[Landauer Cost of Actualization]
\label{thm:landauer}
\textup{[Lean-verified.]}
Each irreversible actualization event that eliminates $\Delta q$ bits of quantum
uncertainty dissipates at minimum
\begin{equation}
  W_{\mathrm{act}} \;=\; \Delta q \cdot k_B T \cdot \ln 2.
  \label{eq:landauer_rahc}
\end{equation}
\end{theorem}

\begin{proof}
By Landauer's principle~\citep{Landauer1961}: the von Neumann entropy decrease
$\Delta S_{\mathrm{vN}} = \Delta q \cdot \ln 2$ of the system must be transferred
to the environment as heat $\geq T \Delta S_{\mathrm{vN}} = k_B T \cdot \Delta q
\cdot \ln 2$.
\end{proof}

\begin{definition}[Steering Cost Functional]
\label{def:steering}
For a forced actualization sequence $\pi = (v_1, \ldots, v_n)$:
\begin{align}
  W_{\mathrm{steer}}(\pi)
  &\;:=\;
  \sum_{k=1}^{n}
  \Bigl[ W_{\mathrm{act}}(v_k) + W_{\mathrm{sel}}(v_k) \Bigr],
  \label{eq:Wsteer_rahc} \\
  W_{\mathrm{sel}}(v_k)
  &\;=\;
  k_B T \cdot \ln\!\left(
    \frac{|\AP_{\mathrm{free}}(\GG_{k-1})|}{|\AP_{\mathrm{forced}}(\GG_{k-1})|}
  \right).
  \label{eq:Wsel_rahc}
\end{align}
\end{definition}

\begin{theorem}[Epistemic-Physical Equivalence]
\label{thm:steering}
\textup{[Lean-verified bound.]}
\begin{equation}
  W_{\mathrm{steer}}(\pi) \;\geq\; k_B T \cdot \Delta F[q],
  \label{eq:EPE_rahc}
\end{equation}
where $\Delta F[q]$ is the reduction in variational free energy (in nats) achieved
by $\pi$.
\end{theorem}

\begin{proof}
$\Delta F[q] = \sum_k \Delta q_k$ (nats). Cost per step $\geq k_B T \cdot \Delta
q_k \cdot \ln 2$ by Theorem~\ref{thm:landauer}; $W_{\mathrm{sel}} \geq 0$.
Summing: $W_{\mathrm{steer}} \geq k_B T \cdot \Delta F[q]$.
\end{proof}

The Action Value function of Section~\ref{sec:steering} is built directly on
Theorems~\ref{thm:landauer} and~\ref{thm:steering}. The biological and
evolutionary applications of these bounds are developed in
RACI~\citep{Sandeman2026c}.

% ─────────────────────────────────────────────────────────────────────────────
\section{The Mathematics of Graph Steering}
\label{sec:steering}
% ─────────────────────────────────────────────────────────────────────────────

\subsection{Setup: Agency as Causal Forcing}

This section derives the Action Value function for agentic graph steering directly
from the Steering Cost Functional of RACI~\citep{Sandeman2026c}. No additional
assumptions are introduced beyond those already established in the RA suite.

Under passive RA-CSG dynamics, the adjacent possible
$\AP_{\mathrm{free}}(\GG_n)$ at stage $n$ contains all on-shell transitions
available from the current graph configuration. An agent with a Tier 4 internal
model intervenes by restricting the adjacent possible to a forced subset:
\begin{equation}
  \AP_{\mathrm{forced}}(\GG_n) \;\subsetneq\; \AP_{\mathrm{free}}(\GG_n).
\end{equation}
By Definition~\ref{def:steering} (Steering Cost Functional,
Section~\ref{sec:steering_rahc}), the thermodynamic cost of a forced actualization
sequence $\pi = (v_1, \ldots, v_k)$ is:
\begin{equation}
  W_{\mathrm{steer}}(\pi)
  \;:=\;
  \sum_{j=1}^{k}
  \Bigl[
    W_{\mathrm{act}}(v_j) + W_{\mathrm{sel}}(v_j)
  \Bigr],
  \label{eq:Wsteer}
\end{equation}
where $W_{\mathrm{act}}(v_j) = \Delta q_j \cdot k_BT \ln 2$ is the Landauer cost
(Theorem~\ref{thm:landauer}) and $W_{\mathrm{sel}}(v_j)$ is the selection cost
(equation~\eqref{eq:Wsel_rahc}).

By Theorem~\ref{thm:steering} (Epistemic-Physical Equivalence,
Section~\ref{sec:steering_rahc}):
\begin{equation}
  W_{\mathrm{steer}}(\pi) \;\geq\; k_BT \cdot \Delta F[q],
  \label{eq:steer_bound}
\end{equation}
where $\Delta F[q]$ is the reduction in variational free energy (in nats) achieved
by executing $\pi$. The factor $k_BT$ converts nats to joules.

\subsection{The Three Components of Action Value}

An agentic system selects the forced actualization sequence $\pi^*$ that
maximises a balance between expected epistemic gain, pragmatic benefit, and
thermodynamic cost. We formalize this as the \emph{Action Value} of a candidate
sequence $\pi$ given the current graph state $\GG_n$:

\begin{definition}[Action Value Function]
\label{def:actionvalue}
For a candidate forced actualization sequence $\pi$ executed from graph state
$\GG_n$, the Action Value is:
\begin{equation}
  V(\pi \mid \GG_n)
  \;=\;
  \alpha \cdot I(\pi;\, \GG_{n+1})
  \;+\;
  \beta \cdot \mathbb{E}\bigl[U(\GG_{n+1}) \mid \pi,\, \GG_n\bigr]
  \;-\;
  \gamma \cdot W_{\mathrm{steer}}(\pi),
  \label{eq:actionvalue}
\end{equation}
where $\alpha, \beta, \gamma > 0$ are weighting parameters, and the three terms
are defined below.
\end{definition}

\paragraph{Term 1: Epistemic Forcing ($I$).}
$I(\pi;\, \GG_{n+1})$ is the mutual information between executing $\pi$ and the
resulting graph state $\GG_{n+1}$---the expected reduction in the agent's
uncertainty about the future graph structure that results from executing $\pi$:
\begin{equation}
  I(\pi;\, \GG_{n+1})
  \;=\;
  H(\GG_{n+1}) - \mathbb{E}_{\pi}\bigl[H(\GG_{n+1} \mid \pi)\bigr],
\end{equation}
where $H$ denotes Shannon entropy over the probability distribution on future graph
states induced by the RA-CSG dynamics. This term is high when the agent's internal
model $\GG_{\mathrm{model}}$ has high uncertainty about the future, and low when
the future is already well-characterized.

\emph{RA derivation.} The mutual information $I(\pi;\, \GG_{n+1})$ is the
information-theoretic dual of the relative entropy
$S(\rho \| \sigma_0)$ that serves as the actualization criterion in RAQM
and RAQI~\citep{Sandeman2026a,Sandeman2026d}. Executing $\pi$ forces specific
actualization events that increase $S(\rho \| \sigma_0)$ irreversibly, thereby
reducing the agent's uncertainty about the graph state. The epistemic term is
therefore the information gain from \emph{deliberately triggering} the relative
entropy increase that constitutes actualization---using the Kinematic Snap as an
instrument of perception.

\paragraph{Term 2: Pragmatic Steering ($U$).}
$\mathbb{E}[U(\GG_{n+1}) \mid \pi, \GG_n]$ is the expected utility of the
graph state resulting from $\pi$. The utility function $U$ encodes the agent's
autopoietic objective: maintaining the criticality parameter $\sigmaedge \approx 1$.
Formally:
\begin{equation}
  U(\GG_{n+1})
  \;=\;
  -\bigl(\sigmaedge(\GG_{n+1}) - 1\bigr)^2
  \;-\;
  \kappa_{\mathrm{threat}} \cdot \mathbb{P}[\text{Markov blanket breach} \mid \GG_{n+1}],
  \label{eq:utility}
\end{equation}
where the first term penalises deviation from the edge-of-chaos regime and the
second term penalises configurations that threaten Markov closure. The agent
selects $\pi^*$ that steers $\GG_{n+1}$ toward $\sigmaedge \approx 1$ and away
from Markov blanket instability.

\emph{RA derivation.} The utility function~\eqref{eq:utility} is not an
exogenous preference imposed on the framework; it is the autopoietic objective
derived from the definition of Tier 3 systems in Section~\ref{sec:continuum}.
A system is autopoietic if and only if it maintains $\sigmaedge \approx 1$. The
utility function is therefore the formal expression of the Tier 3 survival
condition, elevated to an explicit optimization target by the Tier 4 agent.

\paragraph{Term 3: Thermodynamic Cost ($W_{\mathrm{steer}}$).}
$W_{\mathrm{steer}}(\pi)$ is the steering cost functional~\eqref{eq:Wsteer},
comprising the Landauer cost of actualization and the selection cost of
restricting the adjacent possible. This term prevents the agent from pursuing
unlimited graph forcing: by the Epistemic-Physical Equivalence
bound~\eqref{eq:steer_bound}, the thermodynamic cost of any forced sequence is
bounded below by $k_BT \cdot \Delta F[q]$. An agent cannot force continuous,
unlimited actualizations because its available free energy budget is finite.

Crucially, the cost term applies to \emph{internal} simulation sequences as well
as external actions. The generative motif (Section~\ref{subsec:generative})
computes predictions of $\GG_{\mathrm{env}}$ via real, on-shell actualization
events within the agent's own graph. Each such internal simulation step is a
forced actualization sequence $\pi_{\mathrm{sim}}$ with cost:
\begin{equation}
  W_{\mathrm{steer}}(\pi_{\mathrm{sim}}) \;\geq\; k_BT \cdot \Delta F[q_{\mathrm{sim}}],
  \label{eq:sim_cost}
\end{equation}
where $\Delta F[q_{\mathrm{sim}}]$ is the variational free energy reduced by
running the simulation. The Action Value function therefore governs both physical
actions (movements, communications, resource acquisition) and cognitive actions
(simulating counterfactuals, updating the generative motif, planning). There is no
free lunch in cognition: every thought that reduces uncertainty about the future
pays the Landauer floor cost for the uncertainty it resolves. This directly
accounts for the high metabolic cost of predictive neural processing.

\begin{remark}[Correspondence with Friston's Expected Free Energy]
The Action Value function~\eqref{eq:actionvalue} has a precise correspondence
with Friston's Expected Free Energy $G(\pi)$~\citep{Friston2010,FristonEtAl2017},
which decomposes as:
\begin{equation}
  G(\pi)
  \;=\;
  \underbrace{-\,I(\pi;\,\GG_{n+1})}_{\text{epistemic value}}
  \;+\;
  \underbrace{\mathbb{E}\bigl[\log q(\GG_{n+1}\mid\pi) - \log \tilde{p}(\GG_{n+1})\bigr]}_{\text{risk (pragmatic cost)}},
  \label{eq:friston_G}
\end{equation}
where $\tilde{p}$ is the agent's prior over preferred states. In Friston's
framework, the agent selects the policy $\pi$ that minimizes $G(\pi)$; in the
ACD framework, the agent maximizes $V(\pi \mid \GG_n)$.

The two are related by:
\begin{equation}
  V(\pi \mid \GG_n)
  \;\approx\;
  -\,G(\pi) \;-\; \gamma\,W_{\mathrm{steer}}(\pi),
  \label{eq:V_G_correspondence}
\end{equation}
where the correspondence is:
\begin{itemize}[noitemsep]
  \item $\alpha\,I(\pi;\,\GG_{n+1}) \leftrightarrow$ epistemic value in $-G(\pi)$:
    the expected reduction in uncertainty about the future graph state.
  \item $\beta\,\mathbb{E}[U(\GG_{n+1})] \leftrightarrow$ pragmatic value
    $-\mathbb{E}[\log q - \log\tilde{p}]$: the expected alignment of the resulting
    state with the agent's preferred states (here, $\sigmaedge \approx 1$ and
    Markov closure).
  \item $\gamma\,W_{\mathrm{steer}}(\pi) \leftrightarrow$ the additional
    thermodynamic cost term, which is absent in the standard Free Energy
    Principle but is mandatory in the RA account: every forced actualization
    sequence---including internal simulation---pays the Landauer floor.
\end{itemize}

The ACD framework therefore provides the \emph{physical substrate} for which
the Free Energy Principle is the \emph{functional description}: Friston's
$G(\pi)$ is the information-theoretic shadow of the thermodynamic reality
captured by $V(\pi \mid \GG_n)$. The RA derivation adds one term that is
invisible in the standard Free Energy Principle but physically non-negotiable:
the Landauer cost $\gamma\,W_{\mathrm{steer}}$, which enforces the
thermodynamic limits on cognitive capacity identified in Proposition~\ref{prop:bound}.
\end{remark}

\subsection{The Optimal Steering Policy}

The agent's optimal policy selects the forced actualization sequence $\pi^*$ that
maximises the Action Value:
\begin{equation}
  \pi^* \;=\; \operatorname{argmax}_{\pi \in \Pi_{\mathrm{admissible}}}
  V(\pi \mid \GG_n),
  \label{eq:optimal}
\end{equation}
where $\Pi_{\mathrm{admissible}}$ is the set of all causal sequences available to
the agent given its current Markov blanket boundary and free energy budget.

Several properties of~\eqref{eq:optimal} follow directly from the RA framework:

\begin{proposition}[Thermodynamic Bound on Steering]
\label{prop:bound}
For any admissible policy $\pi$,
\begin{equation}
  V(\pi \mid \GG_n) \;\leq\;
  \alpha \cdot I(\pi;\, \GG_{n+1})
  + \beta \cdot \mathbb{E}[U(\GG_{n+1})]
  - \gamma \cdot k_BT \cdot \Delta F[q(\pi)].
\end{equation}
The thermodynamic cost sets a lower bound on the free energy that must be
expended to achieve any given epistemic or pragmatic gain.
\end{proposition}

\begin{proof}
Substitute the bound~\eqref{eq:steer_bound} into Definition~\ref{def:actionvalue}.
\end{proof}

\begin{corollary}[Finite Intelligence]
No agent operating within a finite free energy budget can maximize both epistemic
and pragmatic value simultaneously. The Action Value function encodes a
thermodynamically enforced trade-off between knowing more (high $I$) and doing more
(high $U$) that is fundamental to all intelligent systems, not a contingent
architectural limitation.
\end{corollary}

\subsection{Surprise-Triggered Versus Value-Scored Forcing}

Not all forced actualization sequences are selected by continuous optimization
of~\eqref{eq:optimal}. Some forcing events are triggered by threshold crossings
in the agent's internal model. We distinguish two modes:

\begin{definition}[Surprise-Triggered Forcing]
A surprise-triggered forcing event occurs when the relative entropy of the
agent's current state estimate against the expected state exceeds a threshold:
\begin{equation}
  S(\rho_{\mathrm{current}} \| \rho_{\mathrm{expected}}) \;>\; \theta_{\mathrm{force}},
\end{equation}
at which point the agent executes the forced sequence $\pi^*_{\mathrm{threat}}$
that minimizes the probability of Markov blanket breach.
\end{definition}

\begin{definition}[Value-Scored Forcing]
A value-scored forcing event is selected by continuous optimization of the
Action Value function~\eqref{eq:actionvalue} across all admissible sequences,
without a threshold trigger. This mode governs proactive, anticipatory steering
when no immediate threat is detected.
\end{definition}

The two modes correspond to the fight-or-flight versus deliberative planning
distinction in neuroscience, and to the distinction between reactive and
proactive control in artificial systems. Both are instances of the same underlying
optimization~\eqref{eq:optimal}; the surprise trigger is a computationally efficient
approximation that prioritizes immediate threat responses over full policy
evaluation.

% ─────────────────────────────────────────────────────────────────────────────
\section{Empirical Predictions}
\label{sec:predictions}
% ─────────────────────────────────────────────────────────────────────────────

The ACD framework makes several predictions that distinguish it from standard
complex systems theory and are in principle testable.

\paragraph{P1: Causal depth of metabolic networks.}
By the linear scaling result of RACI~\citep{Sandeman2026c} ($\ARA(M) \approx
\kappa \cdot A(M)$ in condensed phases), evolutionary selection pressure should
drive organisms toward biosynthetic pathways of minimal Assembly Index $A$.
ACD predicts that across evolutionary lineages, organisms should exhibit
systematic reduction in the Assembly Index of their critical molecules over
evolutionary time, independently of thermodynamic cost. This is testable by
comparing Assembly Index values across metabolic pathways in organisms at
different evolutionary distances.

\paragraph{P2: Tier 3/4 transition and the internal model threshold.}
ACD predicts a sharp transition in information processing at the Tier 3/4 boundary:
the emergence of a persistent internal model $\GG_{\mathrm{model}}$ should
correspond to a qualitative change in the agent's ability to perform
counterfactual causal reasoning. Systems below the transition can react to
environmental perturbations but cannot anticipate them. Systems above the
transition can simulate counterfactual graph extensions and select actions that
prevent perturbations before they arrive. This transition is testable in
neuroscience via the presence or absence of prospective coding in neural
activity.

\paragraph{P3: Thermodynamic cost of intelligence.}
Proposition~\ref{prop:bound} implies a strict thermodynamic cost for epistemic
gain. ACD predicts that the metabolic cost of intelligent behavior scales with
the mutual information acquired per unit time, with a proportionality constant
set by $k_BT$. This is consistent with known metabolic scaling of neural tissue
but provides a more precise bound: the free energy cost per bit of information
acquired is at least $k_BT \ln 2$, the Landauer minimum.

\paragraph{P4: Optimal coarse-graining scale.}
The recursive hierarchy $(\Pi_n)$ predicts that stable Markov blankets exist only
at specific scales determined by the causal structure of the underlying graph.
ACD predicts that biological and cognitive systems will exhibit characteristic
scales at which information integration is maximized---corresponding to the
scales at which the coarse-graining projection $\Pi_n$ achieves maximum causal
closure. This is testable via integrated information theory
measures~\citep{Tononi2004} applied at multiple spatial and temporal scales.

% ─────────────────────────────────────────────────────────────────────────────
\section{Relationship to the RA Suite}
\label{sec:suite}
% ─────────────────────────────────────────────────────────────────────────────

The present paper extends the RA research program in a specific direction. Its
relationship to the four companion papers is as follows:

\begin{description}[leftmargin=1.5em]
  \item[\textbf{RAQM}~\citep{Sandeman2026a}] Provides the foundational actualization
    criterion ($p^{\mu}p_{\mu} = m^2$, the Kinematic Snap) and the relative
    entropy criterion for actualization events. ACD uses these as the Level 0
    building blocks of the recursive hierarchy.

  \item[\textbf{RAGC}~\citep{Sandeman2026b}] Provides the Local Ledger Condition,
    the RA-CSG growth dynamics, and the adjacent possible formalism. ACD's
    coarse-graining projection $\Pi_n$ is built on top of the CSG dynamics.

  \item[\textbf{RACI}~\citep{Sandeman2026c}] Develops the biological and
    evolutionary applications of the theorems proved in Section~\ref{sec:theorems}
    of the present paper. Specifically, RACI provides: the full C1$'$ variational
    argument and B3LYP/6-311+G* computational survey establishing the Assembly
    Index Correspondence; the detailed Causal Firewall derivation and edge-of-chaos
    phase transition; England compatibility and evolutionary predictions; the
    complete treatment of agency in biological context; and the four empirical
    predictions for astrobiology, evolutionary biology, and origin-of-life
    research. The division of labour is precise: RAHC proves the foundational
    theorems and establishes the universal hierarchy; RACI applies them to the
    specific domain of biological complexity and life.

  \item[\textbf{RAQI}~\citep{Sandeman2026d}] Provides the Kinematic Coherence
    Bound and the connection between actualization timescales and information
    processing limits. ACD's prediction P3 (thermodynamic cost of intelligence)
    is a macroscopic extension of the RAQI bounds on quantum computing systems.

  \item[\textbf{RADM}~\citep{Sandeman2026RADM}] Extends the actualization-density
    field to galactic scales. The critical actualization density $\lambda_c$ at
    which the causal DAG undergoes dimensional reduction from $d_s = 3$ to
    $d_s = 2$ (Durhuus-Jonsson-Wheater theorem) is a cosmological instance of
    the same phase-transition mechanism that ACD identifies as a tier boundary
    in the complexity hierarchy. The sparse-graph regime ($\mu \ll 1$) in RADM
    is structurally analogous to the pre-Causal-Firewall abiotic chemistry
    regime in RACI --- both involve actualization graphs below the threshold
    for stable self-referential structure.

  \item[\textbf{EB}~\citep{Sandeman2026EB}] Defines the Covariant Stochastic
    Graph Rewriting System whose Step~5 coarse-graining underlies $\Pi_n$.
    Recent progress: the BDG sign criterion ($S_{\mathrm{BDG}}(v) > 0$) as the
    discrete actualization filter yields the Poisson-CSG generative measure
    $c_k^{\mathrm{RA}}(x) \propto e^{-\mu}\mu^k$, making $c_k$ local and
    environment-dependent --- a structural advance that closes the measure
    problem for RA's CSG dynamics.
\end{description}

% ─────────────────────────────────────────────────────────────────────────────
\section{Open Problems and Hard Walls}
\label{sec:hardwalls}
% ─────────────────────────────────────────────────────────────────────────────

Several open problems remain before ACD becomes a complete theory.

\paragraph{Hard Wall 1: The Causal Firewall Limit Theorem --- Resolved.}
The Tier~2/3 boundary corresponds to the critical condition
$\sigma \equiv \lambda\,\tau_d\,\ell^3 = 1$ of RACI~\citep{Sandeman2026c}.
This is now proved rather than conjectured, via the identification of $\sigma$
with the Poisson-CSG parameter $\mu = \lambda\,V_{\mathrm{coh}}\,p_{\mathrm{th}}$
(RAGC~\citep{Sandeman2026b}) and the application of the directed
Erd\H{o}s-R\'{e}nyi theorem~\citep{ErdosRenyi1960}: a directed random causal
graph with mean degree $\mu$ has a giant strongly connected component (shared
causal ledger) if and only if $\mu > 1$. The Causal Firewall threshold $\sigma=1$
is exactly this percolation transition. The trace-distance convergence bound
between observer-dependent actualization histories scales as $\epsilon \sim
N^{-1/3}$ at criticality and decays exponentially above it.
\emph{Remaining gap:} deriving the identification $p_{\mathrm{th}} =
\tau_d/\tau_{\mathrm{mol}}$ from the RA-CSG dynamics (a single calculation
in open quantum systems theory).
\textbf{Collaborator target:} open quantum systems, decoherence theory.

\paragraph{Hard Wall 2: Computational Tractability of $\Pi_n$.}
The coarse-graining projection $\Pi_n$ is defined in terms of indistinguishable
probability distributions over future macro-observables (equation~\eqref{eq:equiv}).
Computing this projection for realistic systems is in general computationally
intractable. Identifying efficiently computable approximations to $\Pi_n$ that
preserve the relevant causal structure is an open problem. \textbf{Collaborator
target:} computational complexity theory, causal discovery algorithms.

\paragraph{Hard Wall 3: Quantitative Utility Function.}
The utility function $U(\GG_{n+1})$ in~\eqref{eq:utility} is defined qualitatively
in terms of the criticality parameter $\sigmaedge$ and the Markov blanket breach
probability. Making this quantitative requires explicit models of both quantities
from the RA-CSG dynamics. The Causal Firewall conjecture provides a candidate
expression for the threshold, but the full utility function requires the limit
theorem of Hard Wall 1. \textbf{Collaborator target:} open quantum systems,
mathematical biology.

\paragraph{Hard Wall 4: Connecting $\Pi_n$ to Integrated Information.}
Prediction P4 connects ACD to integrated information theory~\citep{Tononi2004}.
Making this connection rigorous requires showing that the scales at which
$\Pi_n$ achieves maximum causal closure correspond to the scales at which
integrated information $\Phi$ is maximized. This would provide ACD with a
precise operational criterion for identifying the ``natural'' levels of
description of a complex system. \textbf{Collaborator target:} information
theory, consciousness research.

% ─────────────────────────────────────────────────────────────────────────────
\section{Conclusion}
\label{sec:conclusion}
% ─────────────────────────────────────────────────────────────────────────────

Algorithmic Causal Dynamics provides a unified physical account of the hierarchy
of complexity. By grounding every tier---from on-shell boson exchanges to
metabolic cells to reasoning minds---in the same formal mechanism of recursive
causal closure, the conceptual boundaries between fundamental physics, biological
autopoiesis, and cognition dissolve. Not because they are ``really'' the same
thing, but because they are different regimes of the same topological process.

The key contributions of this paper are:
\begin{enumerate}[noitemsep]
  \item A rigorous topological definition of irreducible emergence as
    \emph{causal closure}: a subgraph emerges at scale $n$ when it achieves a
    stable Markov blanket and can be treated as a single effective causal node
    at scale $n+1$. This is the first such definition grounded in causal graph
    physics rather than information-theoretic postulates.
  \item The recursive coarse-graining hierarchy $(\Pi_1, \Pi_2, \ldots, \Pi_n)$,
    which generates the full four-tier continuum of physical complexity---from
    on-shell actualization events through dissipative structures, autopoietic
    systems, and agentic intelligence---as successive regimes of causal closure
    produced by a single formal mechanism.
  \item The Generative Motif definition and the Tier~4 Transition Criterion,
    providing the first precise topological criterion for distinguishing reactive
    autopoiesis from predictive agency in terms of internal causal graph structure.
  \item A derivation of the Action Value function for agentic graph-steering
    directly from the RA thermodynamic framework, establishing the formal
    correspondence $V(\pi|\GG_n) \approx -G(\pi) - \gamma W_{\mathrm{steer}}(\pi)$
    that positions RA as the physical substrate of Friston's Free Energy Principle.
  \item Four empirical predictions that distinguish ACD from standard complex
    systems theory and connect to neuroscience, astrobiology, and integrated
    information theory.
\end{enumerate}

The hierarchy is not a ladder climbing toward intelligence as its goal; it is
a physical structure in which each tier is a stable regime of causal closure, as
real and as fundamental as any other. Chemistry is not ``merely'' physics, nor
is biology ``merely'' chemistry, nor cognition ``merely'' biology. Each is a
distinct topological regime generated by the same underlying mechanism and
irreducible to the tier below it---not because reductionism fails in principle,
but because the macroscopic causal edges that govern behaviour at each tier are
structurally absent at lower scales.

Intelligence, at the top of this hierarchy, is not mysterious and not special
in kind. It is the tier at which causal closure becomes sufficiently deep and
recurrent that a subgraph can model the causal structure of its environment,
anticipate threats to its own Markov closure, and steer the growth of the
universal DAG to ensure its continued actualization. The thermodynamic cost of
this capacity---grounded in Proposition~\ref{prop:bound} and anchored empirically
by the metabolic budget of the brain---is not an evolutionary inefficiency. It is
the Landauer floor for running a physical simulation of a complex world. That the
universe has generated systems willing to pay this cost is one of the most
striking structural facts about the causal graph we inhabit.

% ─────────────────────────────────────────────────────────────────────────────
\begin{thebibliography}{99}
% ─────────────────────────────────────────────────────────────────────────────

\bibitem[Erd\H{o}s \& R\'{e}nyi(1960)]{ErdosRenyi1960}
P.~Erd\H{o}s and A.~R\'{e}nyi.
\newblock On the evolution of random graphs.
\newblock \emph{Publications of the Mathematical Institute of the Hungarian
  Academy of Sciences}, 5:17--61, 1960.

\bibitem[Sandeman(2026e)]{Sandeman2026RADM}
J.~F.~Sandeman.
\newblock Relational Actualism: Dark Matter as Actualization-Density Topology ({RADM}).
\newblock Preprint, 2026.
\newblock \href{https://doi.org/10.5281/zenodo.19198513}{doi.org/10.5281/zenodo.19198513}.

\bibitem[Sandeman(2026g)]{Sandeman2026EB}
J.~F.~Sandeman.
\newblock The Engine of Becoming: A Covariant Stochastic Graph Rewriting Algorithm as the Unified Foundation of Quantum Mechanics and General Relativity ({EB}).
\newblock Preprint, 2026.
\newblock \href{https://doi.org/10.5281/zenodo.19198120}{doi.org/10.5281/zenodo.19198120}.

\bibitem[Sandeman(2026a)]{Sandeman2026a}
J.~F.~Sandeman.
\newblock Relational Actualism and Quantum Mechanics ({RAQM}).
\newblock \textit{Foundations of Physics} (Under review), 2026.
\newblock \href{https://doi.org/10.5281/zenodo.19174380}{doi.org/10.5281/zenodo.19174380}.

\bibitem[Sandeman(2026b)]{Sandeman2026b}
J.~F.~Sandeman.
\newblock Relational Actualism: Gravity and Cosmology ({RAGC}).
\newblock Preprint, 2026.
\newblock \href{https://doi.org/10.5281/zenodo.19197999}{doi.org/10.5281/zenodo.19197999}.

\bibitem[Sandeman(2026c)]{Sandeman2026c}
J.~F.~Sandeman.
\newblock Relational Actualism: Complexity and Intelligence ({RACI}).
\newblock Preprint, 2026.
\newblock \href{https://doi.org/10.5281/zenodo.19198224}{doi.org/10.5281/zenodo.19198224}.

\bibitem[Sandeman(2026d)]{Sandeman2026d}
J.~F.~Sandeman.
\newblock Relational Actualism: Implications for Quantum Information ({RAQI}).
\newblock Preprint, 2026.
\newblock \href{https://doi.org/10.5281/zenodo.19198285}{doi.org/10.5281/zenodo.19198285}.

\bibitem[Attwell and Laughlin(2001)]{Attwell2001}
D.~Attwell and S.~B.~Laughlin.
\newblock An energy budget for signaling in the grey matter of the brain.
\newblock \textit{Journal of Cerebral Blood Flow \& Metabolism}, 21(10):1133--1145, 2001.

\bibitem[Friston et~al.(2017)]{FristonEtAl2017}
K.~Friston, T.~FitzGerald, F.~Rigoli, P.~Schwartenbeck, and G.~Pezzulo.
\newblock Active inference: A process theory.
\newblock \textit{Neural Computation}, 29(1):1--49, 2017.

\bibitem[Friston(2010)]{Friston2010}
K.~Friston.
\newblock The free-energy principle: a unified brain theory?
\newblock \textit{Nature Reviews Neuroscience}, 11(2):127--138, 2010.

\bibitem[Kauffman(2000)]{Kauffman2000}
S.~A.~Kauffman.
\newblock \textit{Investigations}.
\newblock Oxford University Press, Oxford, 2000.

\bibitem[Landauer(1961)]{Landauer1961}
R.~Landauer.
\newblock Irreversibility and heat generation in the computing process.
\newblock \textit{IBM Journal of Research and Development}, 5(3):183--191, 1961.

\bibitem[Maturana and Varela(1980)]{Maturana1980}
H.~R.~Maturana and F.~J.~Varela.
\newblock \textit{Autopoiesis and Cognition: The Realization of the Living}.
\newblock D.~Reidel, Dordrecht, 1980.

\bibitem[Murray et~al.(2023)]{Murray2023}
J.~Murray et al.
\newblock Assembly theory explains and quantifies selection and evolution.
\newblock \textit{Nature}, 622:321--328, 2023.

\bibitem[Pearl(1988)]{Pearl1988}
J.~Pearl.
\newblock \textit{Probabilistic Reasoning in Intelligent Systems}.
\newblock Morgan Kaufmann, San Francisco, 1988.

\bibitem[Tononi(2004)]{Tononi2004}
G.~Tononi.
\newblock An information integration theory of consciousness.
\newblock \textit{BMC Neuroscience}, 5:42, 2004.

\end{thebibliography}

\section*{Declarations}
\textbf{Funding:} Not applicable.\quad
\textbf{Competing interests:} The author declares no competing interests.

\textbf{Preprint availability:} This manuscript is available as a preprint at
\href{https://doi.org/10.5281/zenodo.19198171}{doi.org/10.5281/zenodo.19198171}.

\textbf{Code availability:} All companion paper source files and Lean~4 proof files are publicly available at \href{https://github.com/jsandeman/Relational-Actualism}{github.com/jsandeman/Relational-Actualism}.

\textbf{Data availability:} This paper presents theoretical work. No datasets
were generated or analysed.

\end{document}
